\label{subsubsec:funzioni_dataset}
\subsubsection{get\_all\_datasets}
\paragraph{Descrizione}
La funzione \texttt{get\_all\_datasets} permette di recuperare tutti i dataset salvati nel sistema che contengono la query nel proprio nome.
\paragraph{Dettagli dell'endpoint}
\begin{itemize}
    \item \textbf{HTTP method}: GET;
    \item \textbf{Endpoint}: \texttt{/datasets};
    \item \textbf{Response Model}: \texttt{DatasetsListDTO};
    \item \textbf{Dependency Injection}: 
    \begin{itemize}
    \item \textbf{dataset\_service (DatasetUseCase)}: la dipendenza viene risolta tramite \textit{Dependency Injector}. Questo permette di ottenere l'istanza del \texttt{DatasetService} dal \texttt{Container}, che si occupa di tutte le configurazioni e inizializzazioni necessarie per il servizio.
    \end{itemize}
\end{itemize}
\paragraph{Parametri}
\begin{itemize}
    \item \textbf{q (string: "")}: query per filtrare i dataset salvati.
\end{itemize}
\paragraph{Implementazione}
\begin{itemize}
    \item Effettua una chiamata al servizio di gestione dei dataset per recuperare l'elenco completo dei dataset salvati;
    \item Restituisce lo stato e un oggetto \texttt{DatasetListDTO} nel caso in cui l'operazione sia andata a buon fine.
\end{itemize}

\subsubsection{get\_dataset}
\paragraph{Descrizione}
La funzione \texttt{get\_dataset} permette di recuperare i metadati di un dataset indicato.
\paragraph{Dettagli dell'endpoint}
\begin{itemize}
    \item \textbf{HTTP method}: GET;
    \item \textbf{Endpoint}: \texttt{/datasets/<dataset\_id>};
    \item \textbf{Response Model}: \texttt{DatasetDTO};
    \item \textbf{Dependency Injection}: 
    \begin{itemize}
        \item \textbf{dataset\_service (DatasetUseCase)}: la dipendenza viene risolta tramite \textit{Dependency Injector}. Questo permette di ottenere l'istanza del \texttt{DatasetService} dal \texttt{Container}, che si occupa di tutte le configurazioni e inizializzazioni necessarie per il servizio;
    \end{itemize}
\end{itemize}
\paragraph{Parametri}
\begin{itemize}
    \item \textbf{dataset\_id (int)}: ID del dataset da recuperare.
\end{itemize}
\paragraph{Implementazione}
\begin{itemize}
    \item Effettua una chiamata al servizio di gestione dei dataset per recuperare i dettagli del dataset salvato specificato;
    \item Restituisce lo stato e un oggetto \texttt{DatasetDTO} nel caso in cui l'operazione sia andata a buon fine.
\end{itemize}

\subsubsection{create\_dataset}
\paragraph{Descrizione}
La funzione \texttt{create\_dataset} permette di creare un nuovo dataset nel sistema, che può essere vuoto, popolato a partire da un file JSON o una copia di un dataset già salvato.
\paragraph{Dettagli dell'endpoint}
\begin{itemize}
    \item \textbf{HTTP method}: POST;
    \item \textbf{Endpoint}: \texttt{/datasets};
    \item \textbf{Response Model}: \texttt{DatasetDTO};
    \item \textbf{Dependency Injection}: 
    \begin{itemize}
    \item \textbf{dataset\_service (DatasetUseCase)}: la dipendenza viene risolta tramite \textit{Dependency Injector}. Questo permette di ottenere l'istanza del \texttt{DatasetService} dal \texttt{Container}, che si occupa di tutte le configurazioni e inizializzazioni necessarie per il servizio.
    \end{itemize}
\end{itemize}
\paragraph{Parametri}
\begin{itemize}
    \item \textbf{dataset (DatasetDTO)}: metadati del dataset;
    \item \textbf{file}: file JSON che definisce gli elementi del dataset;
    \item \textbf{copy (int)}: ID del dataset da copiare.
\end{itemize}
\paragraph{Implementazione}
\begin{itemize}
    \item Effettua una chiamata al servizio di gestione dei dataset per creare un nuovo dataset;
    \item Restituisce lo stato e un oggetto \texttt{DatasetDTO} nel caso in cui l'operazione sia andata a buon fine.
\end{itemize}

\subsubsection{update\_dataset}
\paragraph{Descrizione}
La funzione \texttt{update\_dataset} permette di salvare le modifiche apportate a un dataset esistente o appena creato.
\paragraph{Dettagli dell'endpoint}
\begin{itemize}
    \item \textbf{HTTP method}: PUT;
    \item \textbf{Endpoint}: \texttt{/datasets/<dataset\_id>};
    \item \textbf{Response Model}: \texttt{DatasetDTO};
    \item \textbf{Dependency Injection}: 
    \begin{itemize}
        \item \textbf{dataset\_service (DatasetUseCase)}: la dipendenza viene risolta tramite \textit{Dependency Injector}. Questo permette di ottenere l'istanza del \texttt{DatasetService} dal \texttt{Container}, che si occupa di tutte le configurazioni e inizializzazioni necessarie per il servizio;
    \end{itemize}
\end{itemize}
\paragraph{Parametri}
\begin{itemize}
    \item \textbf{dataset\_id (int)}: ID del dataset da aggiornare;
    \item \textbf{dataset (DatasetDTO)}: metadati del dataset.
\end{itemize}
\paragraph{Implementazione}
\begin{itemize}
    \item Effettua una chiamata al servizio di gestione dei dataset per aggiornare il dataset;
    \item Restituisce lo stato e un oggetto \texttt{DatasetDTO} nel caso in cui l'operazione sia andata a buon fine.
\end{itemize}

\subsubsection{delete\_dataset}
\paragraph{Descrizione}
La funzione \texttt{delete\_dataset} permette di eliminare un dataset.
\paragraph{Dettagli dell'endpoint}
\begin{itemize}
    \item \textbf{HTTP method}: DELETE;
    \item \textbf{Endpoint}: \texttt{/datasets/<dataset\_id>};
    \item \textbf{Dependency Injection}: 
    \begin{itemize}
        \item \textbf{dataset\_service (DatasetUseCase)}: la dipendenza viene risolta tramite \textit{Dependency Injector}. Questo permette di ottenere l'istanza del \texttt{DatasetService} dal \texttt{Container}, che si occupa di tutte le configurazioni e inizializzazioni necessarie per il servizio;
    \end{itemize}
\end{itemize}

\paragraph{Parametri}
\begin{itemize}
    \item \textbf{dataset\_id (int)}: ID del dataset da eliminare.
\end{itemize}

\paragraph{Implementazione}
\begin{itemize}
    \item Effettua una chiamata al servizio di gestione dei dataset per eliminare il dataset specificato;
    \item Restituisce l'esisto dell'operazione.
\end{itemize}

\subsubsection{get\_qa\_page}
\paragraph{Descrizione}
La funzione \texttt{get\_qa\_page} permette di recuperare la pagina indicata di un dataset fornito.
Permette anche di filtrare gli elementi del dataset che contengono una data query.

\paragraph{Dettagli dell'endpoint}
\begin{itemize}
    \item \textbf{HTTP method}: GET;
    \item \textbf{Endpoint}: \texttt{/datasets/<dataset\_id>/qas};
    \item \textbf{Response Model}: \texttt{DatasetPageDTO};
    \item \textbf{Dependency Injection}: 
    \begin{itemize}
        \item \textbf{qa\_service (QaUseCase)}: la dipendenza viene risolta tramite \textit{Dependency Injector}. Questo permette di ottenere l'istanza del \texttt{QaService} dal \texttt{Container}, che si occupa di tutte le configurazioni e inizializzazioni necessarie per il servizio;
    \end{itemize}
\end{itemize}
\paragraph{Parametri}
\begin{itemize}
    \item \textbf{dataset\_id (int)}: ID del dataset da recuperare;
    \item \textbf{page (int: 1)}: numero di pagina del dataset;
    \item \textbf{q (string: "")}: query da utilizzare per filtrare gli elementi del dataset.
\end{itemize}
\paragraph{Implementazione}
\begin{itemize}
    \item Effettua una chiamata al servizio di gestione degli elementi per recuperare le domande-risposte che dopo essere state filtrate appartengono alla pagina richiesta;
    \item Restituisce lo stato e un oggetto \texttt{DatasetPageDTO} nel caso in cui l'operazione sia andata a buon fine.
\end{itemize}

\subsubsection{create\_qa}
La funzione \texttt{create\_qa} permette di creare una nuova coppia domanda e risposta per il dataset indicato.
\paragraph{Dettagli dell'endpoint}
\begin{itemize}
    \item \textbf{HTTP method}: POST;
    \item \textbf{Endpoint}: \texttt{/datasets/<dataset\_id>/qas};
    \item \textbf{Response Model}: \texttt{QaDTO};
    \item \textbf{Dependency Injection}: 
    \begin{itemize}
        \item \textbf{qa\_service (QaUseCase)}: la dipendenza viene risolta tramite \textit{Dependency Injector}. Questo permette di ottenere l'istanza del \texttt{QaService} dal \texttt{Container}, che si occupa di tutte le configurazioni e inizializzazioni necessarie per il servizio;
    \end{itemize}
\end{itemize}

\paragraph{Parametri}
\begin{itemize}
    \item \textbf{dataset\_id (int)}: ID del dataset a cui deve essere associata la coppia domanda e risposta;
    \item \textbf{QaDTO}: rappresentazione dell'elemento da inserire nel dataset.
\end{itemize}

\paragraph{Implementazione}
\begin{itemize}
    \item Effettua una chiamata al servizio di gestione delle coppie domanda e risposta per crearne una nuova;
    \item Restituisce lo stato e un oggetto \texttt{QaDTO} nel caso in cui l'operazione sia andata a buon fine.
\end{itemize}

\subsubsection{update\_qa}
La funzione \texttt{update\_qa} permette di salvare le modifiche apportare a una coppia domanda e risposta appartenente a un dataset indicato.
\paragraph{Dettagli dell'endpoint}
\begin{itemize}
    \item \textbf{HTTP method}: PUT;
    \item \textbf{Endpoint}: \texttt{/datasets/<dataset\_id>/qas/<qa\_id>};
    \item \textbf{Response Model}: \texttt{QaDTO};
    \item \textbf{Dependency Injection}: 
    \begin{itemize}
        \item \textbf{qa\_service (QaUseCase)}: la dipendenza viene risolta tramite \textit{Dependency Injector}. Questo permette di ottenere l'istanza del \texttt{QaService} dal \texttt{Container}, che si occupa di tutte le configurazioni e inizializzazioni necessarie per il servizio;
    \end{itemize}
\end{itemize}

\paragraph{Parametri}
\begin{itemize}
    \item \textbf{dataset\_id (int)}: ID del dataset a cui appartiene la coppia domanda e risposta;
    \item \textbf{qa\_id (int)}: ID della coppia domanda e risposta da aggiornare.
\end{itemize}
\paragraph{Implementazione}
\begin{itemize}
    \item Effettua una chiamata al servizio di gestione delle coppie domanda e risposta per aggiornare una coppia domanda e risposta;
    \item Restituisce lo stato e un oggetto \texttt{QaDTO} nel caso in cui l'operazione sia andata a buon fine.
\end{itemize}

\subsubsection{delete\_qa}
La funzione \texttt{delete\_qa} permette di eliminare una coppia domanda e risposta appartenente al dataset indicato.
\paragraph{Dettagli dell'endpoint}
\begin{itemize}
    \item \textbf{HTTP method}: DELETE;
    \item \textbf{Endpoint}: \texttt{/datasets/<dataset\_id>/qas/<qa\_id>};
    \item \textbf{Dependency Injection}: 
    \begin{itemize}
        \item \textbf{qa\_service (QaUseCase)}: la dipendenza viene risolta tramite \textit{Dependency Injector}. Questo permette di ottenere l'istanza del \texttt{QaService} dal \texttt{Container}, che si occupa di tutte le configurazioni e inizializzazioni necessarie per il servizio.
    \end{itemize}
\end{itemize}

\paragraph{Parametri}
\begin{itemize}
    \item \textbf{dataset\_id (int)}: ID del dataset a cui appartiene la coppia domanda e risposta;
    \item \textbf{qa\_id (int)}: ID della coppia domanda e risposta da eliminare.
\end{itemize}
\paragraph{Implementazione}
\begin{itemize}
    \item Effettua una chiamata al servizio di gestione delle coppie domanda e risposta per eliminare una coppia domanda e risposta;
    \item Restituisce lo stato dell'operazione.
\end{itemize}

\paragraph{Tracciamento dei requisiti}
\begin{tabularx}{\textwidth}{|X|c|}
    \hline
    \textbf{ID} & \textbf{Funzione} \\
    \hline
        RFO-16.05, RFO-16.06, RFO-16.07 & \texttt{get\_all\_datasets}\\
    \hline
       RFO-14.01, RFO-14.02, RFO-14.04, RFO-14.09, RFO-17.01, RFO-17.02 & \texttt{get\_dataset}\\
    \hline
        RFO-17.04, RFO-17.05 & \texttt{create\_dataset}\\
    \hline
        RFO-16.01, RFO-16.02, RFO-16.04, RFO-17.03, RFO-17.04 & \texttt{update\_dataset}\\
    \hline
        RFO-16.03, RFF-19.01, RFF-19.02 & \texttt{delete\_dataset}\\
    \hline
        RFO-14.10, RFO-14.12, RFO-14.14, RFO-15.03, RFO-15.04 & \texttt{get\_qa\_page}\\
    \hline
        RFO-14.05, RFO-15.01 & \texttt{create\_qa}\\
    \hline
        RFO-14.08, RFO-15.01, RFO-15.05, RFO-15.06 & \texttt{update\_qa}\\
    \hline
        RFO-14.07 & \texttt{delete\_qa}\\
    \hline
\end{tabularx}